\chapter{Conclusion and Discussion} \label{chap:conclusion}
\chaptersummary{The final chapter brings the methodological and applied parts of the thesis back together. The overall picture is that strong surrogates can support more label-efficient estimation and planning, but the size and interpretation of those gains depend on surrogate quality and on whether independence assumptions are credible in the protein setting.}

\section{Overview}

The thesis brings together two methodological foundations and one applied extension. The first is that noisy or surrogate labels can be corrected efficiently with a small gold-standard calibration set \parencite{llmjudge_eif}. The second is that prediction quality can be translated into labeled-sample planning through prediction-powered power analysis \parencite{ppi_power}. The applied extension is that these ideas remain useful in protein--protein interaction prediction, provided that the predictive and inferential components are kept distinct and linked through a stable prediction-to-inference handoff.

In the protein application, the final surrogate is a sequence-plus-text ESM2-3B attention ensemble. This model provides the strongest protein-edge surrogate in the study and supports stable prediction-powered summaries on both the internal Intra benchmark and the external BIOGRID-human test split. The main conclusion is not that prediction replaces statistical correction, but that a strong surrogate can make correction and planning substantially more label-efficient when the inferential target is kept explicit.

\section{Main Findings}

Three conclusions are most important.
\begin{enumerate}
    \item \textbf{Strong surrogates can be built for protein-edge prediction.} The best-performing model in this study is not a single classical classifier but a merged-pool neural attention ensemble using both sequence and text information.
    \item \textbf{Prediction-powered inference works well once the surrogate is strong.} For the final model carried into the inference layer, prediction-powered inference and \texttt{PPI++} estimates on the main test splits remain close to the full-data references even when the calibration budget is limited.
    \item \textbf{Independence assumptions matter in graph-structured protein-protein interaction data.} The published i.i.d. planning formulas are still useful, but their apparent sample-size savings should be interpreted as baseline results rather than dependence-robust guarantees.
\end{enumerate}

A complementary negative result is also informative. A supplementary tuning campaign, including a log-loss-focused ensemble comparison as well as alternative losses, early-stopping criteria, reduced hard-negative pressure, and post-hoc calibration, did not produce a replacement surrogate that beat the frozen attention ensemble without sacrificing ranking performance. This suggests that, in the present setting, representation quality and evaluation design matter more than small head-level tuning changes.

\section{Limitations}

Several limitations remain. First, prediction-powered inference and planning are applied here only through their published i.i.d.\ formulations. Because protein--protein interaction data are naturally graph-structured, shared-endpoint dependence across observed edges may affect the variance calculations and sample-size formulas used in this chapter. Second, the planning analysis is restricted to mean- and proportion-type estimands, such as prevalence and related classification summaries, rather than to direct differences in predictive metrics such as PR-AUC or ROC-AUC. Third, the external robustness results still depend on the way the BIOGRID benchmark is mapped, filtered, and paired with negative examples. Fourth, in the protein application the surrogate remains only moderately informative for the inferential targets studied: the pilot \(\rho_{Y\hat p}^2\) values are small, so the gains from prediction-powered planning and tuned \texttt{PPI++} are correspondingly limited. Fifth, the predictive pipeline is based on sequence and UniProt text representations and may miss structural, mechanistic, or context-specific biological information that could improve surrogate quality for downstream inference. Finally, the strongest surrogate examined in this thesis is computationally expensive, which limits the scope for broader hyperparameter exploration and repeated resampling experiments.


\section{Future Directions}

Several directions would strengthen this line of work. On the methodological side, an important next step is to extend prediction-powered inference and planning beyond the current i.i.d.\ setting by developing asymptotic theory, variance estimators, and sample-size formulas adapted to graph-dependent edge observations. It would also be useful to broaden the planning layer beyond mean- and proportion-type targets to include other functionals more directly tied to predictive evaluation. On the applied side, the protein analysis could be strengthened through broader external validation, more systematic and biologically informed negative-edge construction, and improved surrogate models with stronger calibration and richer biological representations, including structural information where available.


\section{Closing Remarks}

At a practical level, this thesis develops a reproducible workflow linking protein-edge prediction, a standardized surrogate handoff, prediction-powered inference, and study planning. At a conceptual level, the main message is equally important: strong machine-learning predictions do not eliminate the need for statistical calibration, but when used carefully they can support stable corrected estimation and reduce gold-standard labeling requirements. At the same time, those gains should be interpreted within the limits of the current theory, especially when prediction-powered methods are transported into graph-structured biological settings.
